\chapter*{Ж}
\addcontentsline{toc}{chapter}{Ж}

\textbf{Жандарка}

50°24'45.9"N 30°37'09.8"E

Озеро, при котором, по улице Днепровской, сохранился уголок частного сектора села Позняки. Жандаркой озеро слыло по меньшей мере с 1950-х.

Судя по старым картам, Жандарка это часть некогда большой речки Тельбин (Толбин). Современное русло Жандарки шире, чем было допустим еще в 1930-х, когда озеро выглядело как причудливо извивающийся среди Позняков и Ново-Позняков дракон.

Фактически, современные Русановский канал и озера Тельбин, Королёк (изначальный Нижний Тельбин) и Жандарка лежат по ходу той былой речки Тельбин. А современный (подчеркиваю что современный) Нижний Тельбин – это рукотворное образование на южной ветке речки Дарницы, пущенной там вне природного русла.

На 1930-е годы Жандарка составляла единое целое с тем давним Нижним Тельбином, разрыв произошел позже. Южной же частью русло, часть коего сопоставима ныне с Жандаркой, доходило до озера Белое, которое лежало там, где ныне 10-й Микрорайон Позняков, а именно гимназия 62, школа 314 и детсад номер 5, вот они находятся на месте озера Белого образца стыка 19-20 годов.

Современная Жандарка до этих месте не дотягивает, она хоть стала шире, но значительно меньше по протяженности.